% =====================================================================
%  c10 —— 第 25–30 节（p037–p042）
%  节标题起始页：25→p037，26→p037，27→p038，28→p039，29→p040，30→p040
%  编号公式（按页面图像实际条数）：63 条
% =====================================================================

% [Tutor Note: 原文 p036 末尾为公式 (24.7)，p037 顶端以接续词「かつ」开头、后接公式 (24.8)–(24.13)，句末「となる」在本页；此处按跨页断句处理，正文从此处直接接续第 24 节。]
且
\begin{align}
  g\left(\frac{\partial}{\partial \tilde y_i}, \frac{\partial}{\partial \tilde y_j}\right)
  &= g\left(\sum_{k=1}^{n}\frac{\partial y_k}{\partial \tilde y_i}\frac{\partial}{\partial y_k},
    \sum_{l=1}^{n}\frac{\partial y_l}{\partial \tilde y_i}\frac{\partial}{\partial y_l}\right) \tag{24.8}\\
  &= \sum_{k,l=1}^{n}\frac{\partial y_k}{\partial \tilde y_i}\frac{\partial y_l}{\partial \tilde y_i}
    g\left(\frac{\partial}{\partial y_k}, \frac{\partial}{\partial y_l}\right) \tag{24.9}\\
  &= \sum_{k,l=1}^{n}\frac{\partial y_k}{\partial \tilde y_i}\frac{\partial y_l}{\partial \tilde y_i}
    \frac{\partial \tilde y_k}{\partial y_l} \tag{24.10}\\
  &= \sum_{k=1}^{n}\frac{\partial y_k}{\partial \tilde y_i}\frac{\partial y_k}{\partial \tilde y_j} \tag{24.11}\\
  &= \sum_{k=1}^{n}\frac{\partial y_j}{\partial \tilde y_i} \tag{24.12}\\
  &= \frac{\partial^2 \tilde K}{\partial \tilde y_i\partial \tilde y_j} \tag{24.13}
\end{align}
成立。从而，$(U,g)$ 以 $((\tilde y_i)_{i=1}^{n})$ 为仿射坐标而具有 Hessian 结构；这称为对偶 Hessian 结构（dual Hesse structure）。又，从 $((y_i)_{i=1}^{n},K)$ 到 $((\tilde y_i)_{i=1}^{n},\tilde K)$ 的变换称为 Legendre 变换（Legendre transformation）。这一构造经粘合给出 $M$ 的对偶 Hessian 结构。

% [Tutor Note: 原文 (24.8)–(24.10) 第二个求和因子中偏导的下标均印作 \tilde y_i（按推导似应为 \tilde y_j）；又 (24.12) 印作 \sum_{k=1}^{n}\partial y_j/\partial\tilde y_i，求和指标 k 未在项中出现。以上均按图像照录，请审校。]

\section{辛仿射结构}

设 $\pi: X \to B$ 为 Lagrange 环面纤维丛。对底空间中的点 $b_0$，取其邻域 $U$ 上 $\pi$ 的平凡化 $\pi^{-1}(U) \cong U\times T^n$。令 $\{\gamma_1,\dots,\gamma_n\}$ 为环面 $T^n := \R^n/\Z^n$ 上 1-闭链的标准基，作为 $\omega$ 沿 $\gamma_i$ 的积分而定义
\begin{equation}
  \omega_i := \int_{\gamma_i}\omega \tag{25.1}
\end{equation}
这是对 $\nu \in T_bB$ 给出
\begin{equation}
  \omega_i(\nu) := \int_{\gamma_i}\iota_{\tilde\nu}(\omega) \tag{25.2}
\end{equation}
的底空间上的 1 次微分形式。其中，$\tilde\nu$ 是 $TX|_{L_b}$ 的截面，使得对任意 $x \in L_b := \pi^{-1}(b)$ 有 $\pi_*(\tilde\nu(x)) = \nu$。由 $L_b$ 的 Lagrange 条件
\begin{equation}
  \omega|_{L_b} = 0 \tag{25.3}
\end{equation}
可知 (25.2) 的右端不依赖 $\tilde\nu$ 的取法而确定。又，由 $\omega$ 是闭形式可知 $\omega_i$ 也是闭形式。从而，若 $U$ 取得足够小，则存在 $U$ 上的函数 $y_i$ 满足
\begin{equation}
  \dd y_i = \omega_i \tag{25.4}
\end{equation}
这样的 $y_i$ 在相差热带仿射变换 $\mathrm{GL}_n(\Z)\ltimes\R^n$ 的作用的意义下唯一确定。在 $U\cap U'$ 上，函数 $\{y_i\}_{i=1}^{n}$ 与 $\{y'_i\}_{i=1}^{n}$ 通过 $\mathrm{GL}_n(\Z)\ltimes\R^n$ 的作用互相联系。由此在 $B$ 上确定热带仿射结构，称之为 $B$ 的辛仿射结构（symplectic affine strcutre）。

% [Tutor Note: 原文英文括注印作 “strcutre”（structure 之误），此处照录；下文 (26.2) 后的 “complex tropical affine strcutre” 同。]

\section{复仿射结构}

同样，令 $\{\Gamma_1,\dots,\Gamma_n\}$ 为 $T^n$ 的 $(n-1)$-闭链的标准基。
\begin{equation}
  \lambda_i = -\int_{\Gamma_i}\Im\Omega \tag{26.1}
\end{equation}
如此设定，则 $\dd\lambda_i = 0$，于是存在 $\{\tilde y_i\}_{i=1}^{n}$ 使得
\begin{equation}
  \lambda_i = \dd\tilde y_i \tag{26.2}
\end{equation}
成立。$B$ 上由此确定以该 $\{\tilde y_i\}_{i=1}^{n}$ 为局部坐标的热带仿射结构，称之为 $B$ 的复仿射结构（complex tropical affine strcutre）。

\section{Hitchin 定理}

\begin{theorem}
McLean 度量是 Hessian 度量，且辛仿射结构与复仿射结构互为对偶。
\end{theorem}

\begin{proof}
选取 $\{\gamma_i\}_{i=1}^{n}$ 与 $\{\Gamma_i\}_{i=1}^{n}$ 使得 $\gamma_i\cdot\gamma_j = \delta_{ij}$。
\begin{equation}
  \delta_{ij} = \dd y_i\left(\frac{\partial}{\partial y_j}\right)
  = \omega_i\left(\frac{\partial}{\partial y_j}\right)
  = \int_{\gamma_i}\iota_{\frac{\partial}{\partial y_j}}\omega \tag{27.1}
\end{equation}
所以，
\begin{equation}
  \left[\iota_{\frac{\partial}{\partial y_j}}\omega\right] \in H^1(L_b) \tag{27.2}
\end{equation}
是与 $\gamma_j$ 对偶的元 $\gamma_j^*$。
\begin{align}
  g_{ij} &:= g\left(\frac{\partial}{\partial y_i}, \frac{\partial}{\partial y_j}\right) \tag{27.3}\\
  &= -\int_{L_b}\left(\iota_{\frac{\partial}{\partial y_i}}\omega\right)
    \wedge\left(\iota_{\frac{\partial}{\partial y_j}}\Im\Omega\right) \tag{27.4}\\
  &= -\int_{L_b}\gamma_i^*\wedge\left(\iota_{\frac{\partial}{\partial y_j}}\Im\Omega\right) \tag{27.5}\\
  &= -\int_{\Gamma_i}\iota_{\frac{\partial}{\partial y_j}}\Im\Omega \tag{27.6}\\
  &= \frac{\partial \tilde y_i}{\partial y_j} \tag{27.7}
\end{align}
成立，而 $g_{ij} = g_{ji}$，故
\begin{equation}
  \frac{\partial \tilde y_i}{\partial y_j} = \frac{\partial \tilde y_j}{\partial y_i} \tag{27.8}
\end{equation}
从而存在某个 $K$ 使得
\begin{equation}
  \tilde y_i = \frac{\partial K}{\partial y_i} \tag{27.9}
\end{equation}
可以这样写。这表明 $g$ 是 Hessian 度量、$(y_i)_{i=1}^{n}$ 与 $(\tilde y_i)_{i=1}^{n}$ 是 Legendre 对偶的仿射坐标。
\end{proof}

% [Tutor Note: (27.1) 前的条件原文印作 \gamma_i\cdot\gamma_j = \delta_{ij}（两因子均作 \gamma），按上下文应为 \gamma_i\cdot\Gamma_j；此处照录。]

\begin{example}
在代数环面 $(\C^\times)^n$ 上以
\begin{equation}
  \omega = \frac{1}{2\pi}\sum_{i=1}^{n}\dd r_i\wedge \dd\theta_i \tag{27.10}
\end{equation}
放入平坦度量。其中，令
\begin{equation}
  z_i = \exp(r_i + \sqrt{-1}\theta_i) \tag{27.11}
\end{equation}
若以
\begin{equation}
  \Omega = \frac{1}{(2\pi\sqrt{-1})^{n-1}}\,
  \frac{\dd z_1}{z_1}\wedge\cdots\wedge\frac{\dd z_n}{z_n} \tag{27.12}
\end{equation}
定义全纯体积形式，则容易看出
\begin{equation}
  y_i = \tilde y_i = r_i \tag{27.13}
\end{equation}
成立。
\end{example}

\section{半平坦镜像对称}

设 $B$ 为热带仿射流形，$\pi: X(B) := TB/T_{\Z}B \to B$ 与 $\check\pi: \check X(B) := T^*B/T^*_{\Z}B$ 为对偶环面纤维化。令 $y = (y_1,\dots,y_n)$ 为 $B$ 的热带仿射坐标，$x = (x_1,\dots,x_n)$ 与 $\check x = (\check x_1,\dots,\check x_n)$ 为对应于 $y$ 的 $\pi$ 以及 $\check\pi$ 的纤维方向坐标。由这些坐标，$X(B)$ 的自然复结构由 $J\big(\frac{\partial}{\partial x_i}\big) = \frac{\partial}{\partial y_i}$ 给出。这也可以说成：$X(B)$ 的复坐标由 $z_i = x_i + \sqrt{-1}y_i$ 给出。又，$\check X(B)$ 的辛结构由
\begin{equation}
  \omega = \sum_{i=1}^{n}\dd\check x_i\wedge \dd x_i \tag{28.1}
\end{equation}
给出。

当 $B$ 具有 Hessian 度量、其势由 $K$ 给出时，$X(B)$ 具有以
\begin{equation}
  \omega = q\sqrt{-1}\partial\bar\partial(K\circ\pi)
  = \frac{\sqrt{-1}}{2}\sum_{j,k=1}^{n}
  \frac{\partial^2 K}{\partial y_j\partial y_k}\dd z_j\wedge \dd\bar z_k \tag{28.2}
\end{equation}
为 Kähler 形式的 Kähler 结构。又，$X(B)$ 具有由
\begin{equation}
  \Omega = \dd z_1\wedge\cdots\wedge \dd z_n \tag{28.3}
\end{equation}
给出的全纯体积形式。$\omega^n$ 成为 $\Omega\wedge\overline{\Omega}$ 的常数倍的必要充分条件是 $K$ 满足实 Monge–Ampère 方程
\begin{equation}
  \det\left(\left(\frac{\partial^2 K}{\partial y_j\partial y_k}\right)_{j,k=1}^{n}\right)
  = \text{常数} \tag{28.4}
\end{equation}
此外，在 $\check X(B)$ 中，复坐标为
\begin{equation}
  \check z_j = \check x_j + \sqrt{-1}\frac{\partial K}{\partial y_j} \tag{28.5}
\end{equation}
Kähler 形式由
\begin{equation}
  \omega = \frac{\sqrt{-1}}{2}\sum_{j,k=1}^{n} g^{jk}\dd z_j\wedge \dd\bar z_k,
  \qquad g_{jk} = \frac{\partial^2 K}{\partial y_j\partial y_k} \tag{28.6}
\end{equation}
给出，而 Calabi–Yau 条件 $\omega^n \propto \Omega\wedge\overline{\Omega}$ 再次成为实 Monge–Ampère 方程
\begin{equation}
  \det\left(\left(\frac{\partial^2 K}{\partial y_j\partial y_k}\right)_{j,k=1}^{n}\right)
  = \text{常数} \tag{28.7}
\end{equation}
满足实 Monge–Ampère 方程的 Hessian 流形称为 Monge–Ampère 流形。

\section{Fourier–SYZ 变换}

将 Lagrange 环面纤维化的截面 $s: B \to \check X(B)$ 写为
\begin{equation}
  s(y) = s_1(y)\check x_1 + \cdots + s_n(y)\check x_n \tag{29.1}
\end{equation}
对此，定义 $X(B)$ 上平凡 $U(1)$ 丛的联络
\begin{equation}
  \nabla = \dd + \sqrt{-1}\sum_{i=1}^{n}s_i(y)\dd x_i \tag{29.2}
\end{equation}
该联络的曲率为
\begin{align}
  F &= \dd\left(\sqrt{-1}\sum_{i=1}^{n}s_i(y)\dd x_i\right) \tag{29.3}\\
  &= -\sqrt{-1}\sum_{i,j=1}^{n}\frac{\partial s_i}{\partial y_j}
    \dd x_i\wedge \dd y_j \tag{29.4}
\end{align}
而其 $(2,0)$ 成分为
\begin{equation}
  F^{2,0} = \frac{1}{2}\sum_{i,j=1}^{n}
  \left(\frac{\partial s_i}{\partial y_j} - \frac{\partial s_j}{\partial y_i}\right)
  \dd z_i\wedge \dd z_j \tag{29.5}
\end{equation}
所以，$s$ 成为 Lagrange 截面的必要充分条件是 $\nabla$ 为正则。

\section{非紧近 Calabi–Yau 流形上的特殊拉格朗日环面纤维丛}

Kähler 流形 $(M,g,J,\omega)$ 与全纯体积形式 $\Omega \in H^0\big(\bigwedge^n T^*_M\big)$ 的组 $(M,g,J,\omega,\Omega)$ 被 Joyce 命名为近 Calabi–Yau 流形（almost Calabi–Yau manifold）（见 [Joy05, Section 4.1]）。与 Calabi–Yau 流形的区别在于不要求 (20.4)。从而，近 Calabi–Yau 流形不必是 Ricci 平坦的。本节依照 [Gro01] 解说非紧近 Calabi–Yau 流形上特殊拉格朗日环面纤维丛的构造。另见 [Gol01, Joy02, GHJ03] 等。

设 $m$ 维环面 $T = (S^1)^m$ 具有到 $n$ 维近 Calabi–Yau 流形 $(X,J,\omega,\Omega)$ 的有效 Hamilton 作用。即，$J$、$\omega$ 以及 $\Omega$ 在 $T$ 的作用下不变，且固定 $X$ 的任意点的 $T$ 的元素只有单位元。令该作用的动量映射为 $\mu: X \to \mathfrak{t}^* \cong \R^m$。必要时将 $X$ 换为稠密开子集，可以假定 $T$ 的作用是自由的、$\mu$ 没有临界点。由于 $T$ 的作用保持 $J$，作为 $T$ 作用的基本向量场可以取全纯（holomorphic）的，将其记作 $\cX_1,\dots,\cX_m$。与实的基本向量场 $X_1,\dots,X_m$ 的关系由
\begin{equation}
  X_i = \cX_i + \overline{\cX}_i, \quad i = 1,\dots,m, \tag{30.1}
\end{equation}
给出。

\begin{example}
考虑 $\C$ 的标准 $S^1$ 作用
\begin{equation}
  e^{\sqrt{-1}\theta} : z \mapsto e^{\sqrt{-1}\theta}z \tag{30.2}
\end{equation}
则全纯基本向量场为
\begin{align}
  \cX &= \sqrt{-1}z\frac{\partial}{\partial z} \tag{30.3}\\
  &= \sqrt{-1}(x + \sqrt{-1}y)\frac{1}{2}
    \left(\frac{\partial}{\partial x} - \sqrt{-1}\frac{\partial}{\partial y}\right) \tag{30.4}\\
  &= \frac{1}{2}\left(-y\frac{\partial}{\partial x} + x\frac{\partial}{\partial y}\right)
    + \frac{\sqrt{-1}}{2}\left(x\frac{\partial}{\partial x} + y\frac{\partial}{\partial y}\right) \tag{30.5}
\end{align}
而实的基本向量场为
\begin{equation}
  X = -y\frac{\partial}{\partial x} + x\frac{\partial}{\partial y} \tag{30.6}
\end{equation}
\end{example}

\begin{theorem}
对任意 $\xi \in \mathfrak{t}^*$，
\begin{equation}
  \Omega_{\mathrm{red}} := \iota(X_1,\dots,X_m)\Omega \tag{30.7}
\end{equation}
在辛商
\begin{equation}
  X/\!/_{\xi}T := \mu^{-1}(\xi)/T \tag{30.8}
\end{equation}
上确定 $(n-m)$ 形式，并且 $(X/\!/_{\xi}T, J_{\mathrm{red}}, \omega_{\mathrm{red}},
\Omega_{\mathrm{red}})$ 成为近 Calabi–Yau 流形。
\end{theorem}

\begin{proof}
由辛商的一般理论，$(X/\!/_{\xi}T, J_{\mathrm{red}}, \omega_{\mathrm{red}})$ 是
Kähler 流形，故只需证明 $\Omega_{\mathrm{red}}$ 确定处处不为零的全纯
$(n-m)$ 形式。对 $\bar p \in X/\!/_{\xi}T$ 处 $X/\!/_{\xi}T$ 的切空间的元
$\bar Y_1,\dots,\bar Y_{n-m}$，将其到 $p \in \mu^{-1}(\bar p)$ 处切空间
$T_p\mu^{-1}(\bar p)$ 的提升记为 $Y_1,\dots,Y_{n-m}$。$Y_1,\dots,Y_{n-m}$ 在
相差加上 $X_1,\dots,X_m$ 的线性组合的自由度下唯一确定，而在
$Y_1,\dots,Y_{n-m}$ 上加 $X_1,\dots,X_m$ 的线性组合时，由 $\Omega$ 的反对称性，
\begin{equation}
  \Omega(X_1,\dots,X_m,Y_1,\dots,Y_{n-m}) \tag{30.9}
\end{equation}
保持不变，故 $\Omega_{\mathrm{red}}$ 在辛商 $X/\!/_{\xi}T$ 上定义微分形式。由于
$\Omega$ 是 $(n,0)$ 形式，
\begin{equation}
  \iota(X_1,\dots,X_m)\Omega
  = \iota(\cX_1+\overline{\cX}_1,\dots,\cX_m+\overline{\cX}_m)\Omega
  = \iota(\cX_1,\dots,\cX_m)\Omega \tag{30.10}
\end{equation}
由于 $\Omega$ 是全纯微分形式、$\cX_1,\dots,\cX_m$ 是全纯向量场，(30.10) 右端的
$\iota(\cX_1,\dots,\cX_m)\Omega$ 也是全纯微分形式。由于 $\Omega$ 是处处不为零的
$n$ 形式、$\cX_1,\dots,\cX_m$ 在 $X$ 的任意点处线性无关，故
$\Omega_{\mathrm{red}}$ 也处处不为零。
\end{proof}

\begin{theorem}
设 $(X,J,\omega,\Omega)$ 为 $n$ 维近 Calabi–Yau 流形。又设给定 $m$ 维实环面
$T = (S^1)^m$ 在 $X$ 上的保持 $\Omega$ 的 Hamilton 作用。令 $T$ 作用的
Hamilton 函数为 $\mu: X \to \mathfrak{t}^*$，以 $(\mathfrak{t}^*)^{\mathrm{reg}}$
表示 $\mu$ 的正则值的集合。进而，给定 $(n-m)$ 维流形 $B'$ 与 $T$ 不变的映射
$\pi': X \to B'$，且对任意 $\xi \in (\mathfrak{t}^*)^{\mathrm{reg}}$，$\pi'$
诱导的映射 $\bar\pi': X/\!/_{\xi}T \to B$ 是 $X/\!/_{\xi}T$ 的特殊拉格朗日环面
纤维丛时，映射
\begin{equation}
  \pi := (\mu,\pi'): X \to \mathfrak{t}^*\times B' \tag{30.11}
\end{equation}
给出 $X$ 的特殊拉格朗日环面纤维丛。
\end{theorem}

% [Tutor Note: 定理 30.3 中原文印作 \bar\pi': X//_{\xi}T \to B，按上下文（及下文证明中取 B'）似应为 \to B'；此处照录。]

\begin{proof}
设 $x \in X$ 为 $\mu$ 的正则点，令对应的正则值为 $\xi := \mu(x)$。在辛商
$X/\!/_{\xi}T := \mu^{-1}(\xi)/T$ 中，以 $[x]$ 表示 $x$ 所代表的点，令过 $[x]$
的 $\bar\pi': X/\!/_{\xi}T \to B'$ 的纤维为
$\overline{L}_{b'} := \bar\pi'^{-1}(b')$。其中 $b' := \bar\pi'([x]) = \pi'(x)$。
令 $\bar Y_1,\dots,\bar Y_{n-m}$ 为 $[x]$ 处 $\overline{L}_{b'}$ 的切空间的基，
将其到 $T_x(\mu^{-1}(\xi))$ 的提升记为 $Y_1,\dots,Y_{n-m}$，则
$\{X_1,\dots,X_m,Y_1,\dots,Y_{n-m}\}$ 构成 $T_xL_b$ 的基。其中
$X_1,\dots,X_m$ 是 $T$ 作用的基本向量场。由 $T$ 的交换性，$X_i$ 之间关于
$\omega$ 正交；由动量映射的定义，$X_i$ 与 $Y_j$ 关于 $\omega$ 正交，故若
$\overline{L}_{b'}$ 是拉格朗日子流形，则 $L_{\xi,b}$ 也是拉格朗日子流形。由
$\Omega_{\mathrm{red}}$ 的定义，
\begin{equation}
  \Im\Omega_{\mathrm{red}}(\bar Y_1,\dots,\bar Y_{n-m})
  = \Im\Omega(X_1,\dots,X_m,Y_1,\dots,Y_{n-m}) \tag{30.12}
\end{equation}
所以，由 $\overline{L}_{b'}$ 是特殊拉格朗日子流形可知 $L_{\xi,b}$ 是特殊拉格朗日
子流形。
\end{proof}

% [Tutor Note: 定理 30.3 的证明中原文印作 T_xL_b（上下文应为 T_xL_{\xi,b}）；此处照录。]

作为定理 30.3 的典型应用例，考虑 $\C^n$ 内区域 $U$ 上的 Kähler 形式、借助函数
\begin{equation}
  \varphi: \R^n \to \R \tag{30.13}
\end{equation}
通过
\begin{align}
  \omega &= \frac{\sqrt{-1}}{2}\partial\bar\partial\varphi
    (|z_1|^2,\dots,|z_n|^2) \tag{30.14}\\
  &= \frac{\sqrt{-1}}{2}\partial\left(\sum_{i=1}^{n}\varphi_i z_i\dd\bar z_i\right) \tag{30.15}\\
  &= \frac{\sqrt{-1}}{2}\left(\sum_{i=1}^{n}\varphi_i\dd z_i\wedge \dd\bar z_i
    + \sum_{i,j=1}^{n}\varphi_{ij}z_i\bar z_j\dd z_j\wedge \dd\bar z_i\right) \tag{30.16}
\end{align}
\begin{equation}
  \mbox{} \tag{30.17}
\end{equation}
给出的情形。其中，令
\begin{equation}
  \varphi_i = \frac{\partial\varphi}{\partial x_i}, \qquad
  \varphi_{ij} = \frac{\partial^2\varphi}{\partial x_i\partial x_j} \tag{30.18}
\end{equation}
以上即为所设。

% [Tutor Note: 原文 p042 中编号 (30.17) 独占一行、其左端无任何公式（疑为原文排版时留空的编号行）；此处照录为带编号 (30.17) 的空显示行，请审校确认。]

\begin{corollary}
映射
\begin{equation}
  f: U \to \R^n, \quad
  z \mapsto \left(\varphi_1|z_1|^2 - \varphi_2|z_2|^2, \dots,
  \varphi_1|z_1|^2 - \varphi_n|z_n|^2,\,
  \Im\left(\sqrt{-1}^{n+1}\prod_{i=1}^{n}z_i\right)\right) \tag{30.19}
\end{equation}
的正则值的原像关于 $\omega$ 与 $\Omega = \dd z_1\wedge\cdots\wedge \dd z_n$ 是
特殊拉格朗日子流形。
\end{corollary}

\begin{proof}
Kähler 形式 $\omega$ 关于以
\begin{equation}
  X_i = \sqrt{-1}\left(z_i\frac{\partial}{\partial z_i}
  - \bar z_i\frac{\partial}{\partial\bar z_i}\right), \quad i = 1,\dots,n \tag{30.20}
\end{equation}
为基本向量场的自然 $T^n$ 作用不变。由
\begin{equation}
  \iota_{X_i}\omega = \varphi_i(z_i\dd\bar z_i + \bar z_i\dd z_i)
  + \sum_{j=1}^{n}\varphi_{ij}|z_i|^2(z_j\dd z_j + \bar z_i\dd z_i)
  = \dd(\varphi_j|z_j|^2) \tag{30.21}
\end{equation}
可知 $T^n$ 作用的动量映射由
\begin{equation}
  z \mapsto \left(\varphi_i|z_i|^2\right)_{i=1}^{n} \tag{30.22}
\end{equation}
给出，故保持全纯体积形式 $\Omega$ 的 $(n-1)$ 维子环面
\begin{equation}
  T := \{\alpha = (\alpha_1,\dots,\alpha_n) \in T^n \mid \alpha_1\cdots\alpha_n = 1\}
  \cong T^{n-1} \tag{30.23}
\end{equation}
的作用的动量映射为
\begin{equation}
  \mu: z \mapsto \left(\varphi_1|z_1|^2 - \varphi_2|z_2|^2, \dots,
  \varphi_1|z_1|^2 - \varphi_n|z_n|^2\right) \tag{30.24}
\end{equation}
存在某个 $1 \le i < j \le n$ 使得 $z_i = z_j = 0$ 的 $\C^n$ 的点是 $\mu$ 的临界点。
正则值处的辛商为 $\C$，其坐标为 $\prod_{i=1}^{n}z_i$，全纯体积形式由
\begin{equation}
  \Omega_{\mathrm{red}} := \iota(X_1 - X_2,\dots,X_1 - X_n)
  \dd z_1\wedge\cdots\wedge \dd z_n
  = \sqrt{-1}^{n-1}\dd\left(\prod_{i=1}^{n}z_i\right) \tag{30.25}
\end{equation}
给出。$T$ 不变的映射
\begin{equation}
  \pi': X \to B' := \R, \quad
  z \mapsto \Im\left(\sqrt{-1}^{n+1}\prod_{i=1}^{n}z_i\right) \tag{30.26}
\end{equation}
在辛商上给出特殊拉格朗日环面纤维化。
\end{proof}

% [Tutor Note: 原文 (30.21) 的第二个求和中各项指标混用了 i 与 j（如 |z_i|^2、(z_j dz_j + \bar z_i dz_i) 与右端 d(\varphi_j|z_j|^2)），疑为原文笔误；此处按图像照录，请审校。]

% CHUNK c10 DONE pages p037–p042
